%% ============================================================
%%  chapitre5.tex — الفصل الخامس: المسؤولية القانونية لكل من
%%                  الشركة الأصلية وشركة المناولة
%% ============================================================

\chapter{المسؤولية القانونية لكل من الشركة الأصلية وشركة المناولة}
\label{chap:chapitre5}

%% ─── مقدمة الفصل ──────────────────────────────────────────
\section*{مقدمة الفصل}
\phantomsection\addcontentsline{toc}{section}{مقدمة الفصل}

تكتسي مسألة توزيع المسؤولية القانونية في علاقات الشغل الثلاثية أهمية قصوى من الناحيتين النظرية والعملية؛ ذلك أن إدخال طرف وسيط بين الأجير والمؤسسة المستفيدة من عمله يثير مخاطر تشتت الالتزامات وتنصل كل طرف من تبعات الإخلال بحقوق الأجراء. ومن ثم، فإن نجاعة الحماية القانونية التي تقرها مدونة الشغل تقاس بقدرتها على تعيين المسؤول الحقيقي عن أداء الأجور، والتصريح بالضمان الاجتماعي، وحماية السلامة الجسدية، والتعويض عن حوادث الشغل والإنهاء التعسفي.

ويجيب هذا الفصل بتفصيل عن \textbf{السؤال السادس} من أسئلة المناقشة الخاصة بوضعية المجموعة \textenglish{8}: \textbf{ما المسؤولية القانونية لكل من الشركة الأصلية وشركة المناولة تجاه الأجراء؟}، وذلك عبر أربعة محاور مترابطة:
\begin{itemize}[leftmargin=*, itemsep=4pt]
    \item \textbf{أولاً:} مسؤولية الشركة الأصلية (المستعملة أو المقاول الأصلي)؛
    \item \textbf{ثانياً:} مسؤولية شركة المناولة (أو مقاولة التشغيل المؤقت)؛
    \item \textbf{ثالثاً:} المسؤولية التضامنية والاحتياطية (الأجور، الضمان الاجتماعي، حوادث الشغل)؛
    \item \textbf{رابعاً:} الجزاءات المدنية والجنائية والإدارية المترتبة عن المخالفات.
\end{itemize}

%% ─────────────────────────────────────────────────────────
\section{مسؤولية الشركة الأصلية}

\subsection{المسؤولية المباشرة للمقاولة المستعملة في التشغيل المؤقت}
على الرغم من أن المقاولة المستعملة ليست هي المتعاقد المباشر مع الأجير المؤقت، إلا أن مدونة الشغل وأحكام القانون المدني والاجتماعي تفرض عليها مسؤوليات مباشرة تنشأ عن ممارستها الفعلية لسلطة التوجيه والإشراف داخل فضاء العمل:
\begin{itemize}[leftmargin=*, itemsep=5pt]
    \item \textbf{المسؤولية المباشرة عن الصحة والسلامة وحوادث الشغل (المادة \textenglish{504}):} تُعد المقاولة المستعملة ملزمةً أصالةً باتخاذ جميع التدابير الوقائية والحمائية لضمان سلامة الأجراء المؤقتين أثناء عملهم داخل منشآتها، وتتحمل المسؤولية القانونية المباشرة عن تأمينهم ضد حوادث الشغل والأمراض المهنية وفق مقتضيات القانون رقم \textenglish{18.12} (الظهير الشريف رقم \textenglish{1.14.190} الصادر في \textenglish{29} دجنبر \textenglish{2014}).
    \item \textbf{المسؤولية عن احترام قيود اللجوء والمدد (المواد \textenglish{496}--\textenglish{500}):} تتحمل الشركة المستعملة المسؤولية عن خرق شروط الاستشارة المسبقة لممثلي الأجراء، أو تشغيل الأجراء المؤقتين خارج الحالات الحصرية الأربع، أو في أشغال خطيرة (المادة \textenglish{497})، أو لتعويض مضربين أو مفصولين اقتصادياً (المادة \textenglish{498})، أو إبقاء الأجراء بعد انقضاء المدة القصوى.
    \item \textbf{المسؤولية الجنائية الشخصية للمدير والمسير (المادة \textenglish{548}):} تقرر المادة \textenglish{548} أن كل مشغل أو مدير أو رئيس طبقاً للمادة \textenglish{7} له داخل المؤسسة، وبتفويض من المشغل، الاختصاص والسلطة الكافية لحمل الأجراء على احترام القانون، يكون مسؤولاً جنائياً عن المخالفات المرتكبة، مع بقاء المشغل مسؤولاً مدنياً عن أداء المصاريف والتعويضات المحكوم بها.
\end{itemize}

\subsection{المسؤولية الرقابية والقضائية للمقاول الأصلي في المناولة}
في إطار عقد المقاولة من الباطن، لا يقف دور الشركة الأصلية عند مجرد دفع الفاتورة التجارية لمقاول الباطن، بل تترتب في ذمتها التزامات رقابية ومسؤوليات قانونية:
\begin{enumerate}[leftmargin=*, itemsep=5pt]
    \item \textbf{واجب السهر والمراقبة (المادتان \textenglish{87} و\textenglish{88}):} متى كان المقاول من الباطن غير مقيد في السجل التجاري ولا صاحب أصل تجاري، وجب على المقاول الأصلي السهر على مراعاة ما تنص عليه أحكام الكتاب الثاني من المدونة في شأن الأجراء (المادة \textenglish{87})؛ كما أوجبت المادة \textenglish{88} على المقاول من الباطن تضمين بطاقة الشغل وورقة الأداء البيانات المحددة تنظيمياً.
    \item \textbf{المسؤولية عن حوادث الشغل داخل المؤسسة (المادة \textenglish{89} فقرة أخيرة):} متى نُفذت أشغال المناولة داخل مؤسسات المقاول الأصلي أو ملحقاتها في نطاق المادة \textenglish{89}، يتحمل هذا الأخير المسؤولية عن التعويض عن حوادث الشغل والأمراض المهنية.
    \item \textbf{المسؤولية بوصفها مشغلاً فعلياً عند المناولة المقنعة:} إذا ثبت أن الشركة الأصلية تمارس التبعية القانونية والفنية المباشرة على عمال شركة المناولة وتدمجهم في نشاطها القار—كما في وضعية المجموعة \textenglish{8}—فإن القضاء يسقط الحجاب التعاقدي الصوري ويحملها مسؤولية المشغل الحقيقي كاملة.
\end{enumerate}

%% ─────────────────────────────────────────────────────────
\section{مسؤولية شركة المناولة}

\subsection{المسؤولية الأصلية كمشغل مباشر عن الأجور والعقود والتصريح الاجتماعي}
سواء تعلق الأمر بمقاول من الباطن حقيقي (المادة \textenglish{87}) أو بمقاولة للتشغيل المؤقت مرخص لها (المادة \textenglish{495})، تظل الشركة الوسيطة هي \textbf{المشغل القانوني المباشر} المسؤول أصالةً عن الوفاء بكافة الحقوق العقدية والقانونية للأجراء:
\begin{itemize}[leftmargin=*, itemsep=5pt]
    \item \textbf{إبرام عقود مكتوبة ومشروعة:} تلتزم مقاولة التشغيل المؤقت بتحرير عقد المهمة مكتوباً وتضمينه جميع البيانات الإلزامية في المادة \textenglish{501}، كما يلتزم المقاول من الباطن باحترام الشكلية الكتابية المنصوص عليها في المادة \textenglish{86} وتسليم بطاقة الشغل وورقة الأداء.
    \item \textbf{أداء الأجور والحد الأدنى القانوني (\textenglish{SMIG}):} الالتزام بأداء الأجر المستحق في آجاله القانونية دون أي اقتطاع غير مشروع تحت مسمى أتعاب أو عمولات وساطة (المادتان \textenglish{480} و\textenglish{485})، وباحترام الحد الأدنى القانوني للأجر في القطاعات غير الفلاحية المحدد عن ساعة الشغل في \textenglish{17,92} درهماً إجمالياً ابتداءً من فاتح يناير \textenglish{2026} (المادة \textenglish{1} من المرسوم \textenglish{2.25.983} والمادة \textenglish{356} من المدونة؛ مع عدم خلط الأجر الإجمالي بالساعي أو الصافي، واعتبار رقم \textenglish{3.422,72} درهماً إجمالياً شهرياً مجرد معادلة توضيحية على أساس \textenglish{191} ساعة لا مبلغاً شهرياً جزافياً موحداً).
    \item \textbf{التصريح ووعاء اشتراكات الضمان الاجتماعي:} الالتزام بتسجيل الأجراء لدى الصندوق الوطني للضمان الاجتماعي (\textenglish{CNSS}) وأداء الاشتراكات وفق الظهير الشريف بمثابة قانون رقم \textenglish{1.72.184} (المعدل بالقانون رقم \textenglish{02.24}). وبموجب الفصل \textenglish{19} من هذا الظهير، يتكون وعاء الاشتراك من مجموع الأجور الإجمالية والتعويضات والمنح والمكافآت والمنافع النقدية والعينية المقبوضة (باستثناء تعويضات المصاريف المهنية المبررة والعناصر المعفاة قانونياً بشروط)، ولا يجوز اختزال التصريح الاجتماعي في مبلغ جزافي يعادل الـ \textenglish{SMIG} متى كان الأجر الفعلي أعلى منه.
    \item \textbf{تحمل تبعات إنهاء عقد الشغل:} ممارسة السلطة التأديبية وفق مساطر مدونة الشغل وتحمل تعويضات الفصل والإشعار والضرر متى ثبت تعسفها في الإنهاء.
\end{itemize}

\subsection{المسؤولية عند الإفلاس وسحب الترخيص واللجوء إلى الكفالة}
توقع المشرع حالة عجز المقاولة الوسيطة أو إفلاسها أو توقفها المفاجئ عن أداء مستحقات العمال، وقرر لذلك ضمانتين قانونيتين:
\begin{itemize}[leftmargin=*, itemsep=5pt]
    \item \textbf{استمرار المسؤولية بعد سحب الترخيص (المادة \textenglish{503}):} لا يؤدي سحب الإذن بالممارسة من مقاولة التشغيل المؤقت (المادة \textenglish{487}) إلى براءة ذمتها؛ بل يظل مسيروها ملزمين بالوفاء الكامل بكافة الالتزامات القائمة تجاه الأجراء وتجاه الصندوق الوطني للضمان الاجتماعي.
    \item \textbf{تنفيذ الأحكام على الكفالة المالية (المادتان \textenglish{482} و\textenglish{488}):} في حالة إعسار مقاولة التشغيل المؤقت أو سحب ترخيصها دون الوفاء بالتزاماتها، تأمر المحكمة المختصة باستعمال الكفالة المالية المودعة لدى صندوق الإيداع والتدبير (\textenglish{CDG}) لأداء مستحقات الأجراء وديون \textenglish{CNSS}.
\end{itemize}

%% ─────────────────────────────────────────────────────────
\section{المسؤولية التضامنية (الأجور، الضمان الاجتماعي، حوادث الشغل)}

\subsection{شروط قيام المسؤولية الاحتياطية والدعوى المباشرة (المادتان 89 و90)}
عالج المشرع المغربي حماية أجراء المقاول من الباطن عند تخاذل مشغلهم عبر إقرار نظام المسؤولية الاحتياطية في المادة \textenglish{89}، والمقرونة بالدعوى المباشرة في المادة \textenglish{90}:

\begin{tcolorbox}[colback=red!5, colframe=accentcolor, arc=4pt,
    title={\textbf{نطاق المسؤولية المشتركة والضمانات في المادتين \textenglish{89} و\textenglish{90}}}]
متى تحقق اجتماع الشروط الثلاثة في المقاول من الباطن (\textbf{حالة الإعسار}، \textbf{عدم التقييد في السجل التجاري}، و\textbf{عدم الانخراط في الصندوق الوطني للضمان الاجتماعي})، يصبح المقاول الأصلي مسؤولاً عن الوفاء بالالتزامات الخمسة التالية في حدود المبالغ المترتبة بذمته للمقاول من الباطن:
\begin{enumerate}[leftmargin=*, itemsep=4pt]
    \item \textbf{أداء أجور الأجراء}، مع مراعاة أجل الإشعار المنصوص عليه في المادة \textenglish{91} (\textenglish{60} يوماً).
    \item \textbf{التعويض عن العطلة السنوية} المؤدى عنها.
    \item \textbf{التعويضات عن الفصل} من الشغل.
    \item \textbf{دفع الاشتراكات الواجب أداؤها} إلى الصندوق الوطني للضمان الاجتماعي (\textenglish{CNSS}).
    \item \textbf{دفع الرسم الخاص بالتكوين المهني}.
\end{enumerate}
وتمنح \textbf{المادة \textenglish{90}} كلاً من الأجراء المتضررين و\textenglish{CNSS} حق ممارسة \textbf{دعوى مباشرة} أمام المحكمة ضد المقاول الأصلي لاستخلاص هذه الحقوق.
\end{tcolorbox}

غير أن فعالية هذه الآلية تظل محدودة عملياً متى كان مقاول الباطن مقيداً شكلياً في السجل التجاري أو منخرطاً في \textenglish{CNSS}، أو متى بادر المقاول الأصلي إلى تصفية مستحقات المقاول من الباطن قبل الإنذار؛ ولذلك يتدخل القضاء عند ثبوت الصورية لإقرار مسؤولية تضامنية فعلية وتصحيحية.

\subsection{نطاق التضامن في حوادث الشغل والأمراض المهنية والوقاية الصحية}
يبرز البعد التضامني بأكثر وضوح وقوة في مجال السلامة الجسدية والأخطار المهنية:
\begin{itemize}[leftmargin=*, itemsep=5pt]
    \item \textbf{في المقاولة من الباطن:} تقرر الفقرة الأخيرة من المادة \textenglish{89} أنه، علاوة على الالتزامات الخمسة المذكورة، يكون المقاول الأصلي مسؤولاً عن التعويض عن حوادث الشغل والأمراض المهنية متى أنجزت الأشغال أو الخدمات داخل مؤسساته أو في ملحقاتها؛ كما يسأل بوصفه المشغل الفعلي متى ثبتت الصورية والتبعية المباشرة.
    \item \textbf{في التشغيل المؤقت:} تُحمل المادة \textenglish{504} المقاولة المستعملة المسؤولية المباشرة عن الوقاية وتأمين الأجراء المؤقتين ضد حوادث الشغل والأمراض المهنية وفق القانون رقم \textenglish{18.12}، بينما تظل الوكالة مسؤولة عن الجوانب التعاقدية والتصريح الاجتماعي العام، مما يخلق تكاملاً تضامنياً لحماية الأجير.
\end{itemize}

%% ─────────────────────────────────────────────────────────
\section{الجزاءات المترتبة عن المخالفات}

\subsection{الجزاءات المدنية وإعادة تكييف العقد والتعويضات}
رغم عدم تنصيص باب الشغل المؤقت صراحة على تحويل تلقائي للعقد، فإن القضاء الاجتماعي المغربي يرتب على خرق الضوابط القانونية في التشغيل المؤقت والمناولة جزاءات مدنية صارمة:
\begin{itemize}[leftmargin=*, itemsep=5pt]
    \item \textbf{إعادة تكييف العقد إلى عقد غير محدد المدة (\textenglish{CDI}):} متى تبين للمحكمة أن العقد أبرم خارج حالات المادة \textenglish{496}، أو خلا من البيانات الإلزامية (المواد \textenglish{499}--\textenglish{501})، أو استمر لسنوات في نشاط دائم للمصنع، فإنها تعيد تكييفه باعتباره عقداً دائماً يربط الأجير مباشرة بالمقاولة المستعملة.
    \item \textbf{احتساب الأقدمية والتعويضات عن الفصل التعسفي:} يترتب على اعتبار المقاولة المستعملة مشغلاً أصلياً حق الأجير عند إنهاء خدمته في المطالبة بالتعويض عن الإشعار (المادة \textenglish{43})، والتعويض عن الفصل (المادة \textenglish{53})، والتعويض عن الضرر (المادة \textenglish{41})، وعلاوة الأقدمية، فضلاً عن تسوية وضعه لدى \textenglish{CNSS} بأثر رجعي.
\end{itemize}

\subsection{الجزاءات الزجرية والإدارية والمالية}
تتوزع العقوبات الزجرية والإدارية المقررة قانوناً لردع المخالفات في مدونة الشغل على المستويات التالية:
\begin{itemize}[leftmargin=*, itemsep=5pt]
    \item \textbf{الغرامات المالية في التشغيل المؤقت (المادة \textenglish{505}):} يعاقب القانون على مخالفة أحكام الباب المتعلق بمقاولات التشغيل المؤقت بغرامة مالية تتراوح بين \textbf{\textenglish{2.000} و\textenglish{5.000} درهم}. ويلاحظ أن هذا المبلغ يظل هزيلاً نسبياً مقارنة بحجم أرباح المقاولات الكبرى.
    \item \textbf{غرامات عرقلة تفتيش الشغل (المادة \textenglish{546}):} كل من يحول دون قيام مفتشي الشغل بمهامهم يعاقب بغرامة تتراوح بين \textbf{\textenglish{25.000} و\textenglish{30.000} درهم}، وتضاعف الغرامة في حالة العود.
    \item \textbf{الجزاءات الإدارية والاجتماعية:} سحب الترخيص الإداري من مقاولة التشغيل المؤقت بقرار معلل ودون تعويض من لدن السلطة الحكومية المكلفة بالشغل (المادة \textenglish{487})؛ ويشار في هذا السياق إلى أن المادة \textenglish{506} نصت على مقتضى انتقالي تاريخي أجاز للمحكمة المختصة إغلاق الوكالات القائمة قبل صدور المدونة إن لم تسو وضعيتها داخل أجل \textenglish{6} أشهر من تاريخ النشر بالجريدة الرسمية (\textenglish{8} دجنبر \textenglish{2003})، وذلك إلى جانب ذعائر التأخير والعقوبات المقررة بموجب نظام الضمان الاجتماعي (ظهير \textenglish{1.72.184} وتعديلاته بالقانون \textenglish{02.24}) عند الإخلال بواجب التصريح بالأجراء.
\end{itemize}

%% ─── خلاصة الفصل ──────────────────────────────────────────
\section*{خلاصة الفصل}
\phantomsection\addcontentsline{toc}{section}{خلاصة الفصل}

يتضح أن المشرع المغربي أرسى مسؤوليات متمايزة بين مقاولة التشغيل المؤقت أو مقاول الباطن (كمشغل قانوني مباشر) وبين المقاولة المستعملة أو الأصلية (المسؤولة عن الصحة والسلامة وحوادث الشغل، والضامنة احتياطياً وفق شروط المادة \textenglish{89}). غير أن شروط المادة \textenglish{89} وضعف غرامة المادة \textenglish{505} تفرض في الواقع الميداني تدخل القضاء لتجاوز الشكلية وتحميل الشركة الأصلية مسؤوليتها الكاملة عند المناولة المقنعة. ولإعمال هذه المسؤوليات في الميدان، يبرز الدور المحوري لجهاز تفتيش الشغل، وهو موضوع الفصل السادس.
